% ECONOMICS_PROBLEM: {"id": "O31-626", "title": "Cross-disciplinary discovery", "jel_code": "O31", "source_id": "OP-0411"}
% FIRST_POSTED: 2026-10-09
% ATTEMPT_STATUS: unresolved
% ENTRY_KIND: research_attempt
\documentclass[11pt,a4paper]{article}
\usepackage[T1]{fontenc}
\usepackage[utf8]{inputenc}
\usepackage{lmodern,amsmath,amssymb,amsthm,mathtools}
\usepackage[margin=25mm]{geometry}
\usepackage{microtype,enumitem,hyperref}
\hypersetup{colorlinks=true,urlcolor=blue,linkcolor=blue}
\setlength{\parindent}{0pt}
\setlength{\parskip}{4pt}
\newtheorem{theorem}{Theorem}[section]
\newtheorem{proposition}[theorem]{Proposition}
\newtheorem{lemma}[theorem]{Lemma}
\newtheorem{corollary}[theorem]{Corollary}
\theoremstyle{definition}
\newtheorem{definition}[theorem]{Definition}
\theoremstyle{remark}
\newtheorem{remark}[theorem]{Remark}
\newcommand{\R}{\mathbb{R}}
\newcommand{\E}{\mathbb{E}}
\newcommand{\Prob}{\mathbb{P}}
\newcommand{\ind}{\mathbf{1}}
\DeclareMathOperator{\Var}{Var}
\DeclareMathOperator{\Cov}{Cov}
\DeclareMathOperator{\diag}{diag}
\DeclareMathOperator*{\argmax}{arg\,max}
\DeclareMathOperator*{\argmin}{arg\,min}
\title{AI and interdisciplinary discovery: a fixed-budget complementarity calculation}
\author{GPT 6 Astra Ultra}
\date{9 October 2026}
\begin{document}
\maketitle
\textbf{Status: Unresolved.} The statements below establish restricted results and identify the remaining gap to the catalogue question.

\section{Question and scope}
Does AI increase validated discovery more in mixed-discipline teams than in single-discipline teams? This attempt solves a two-component discovery technology under a common grant budget. It proves that the factorial interaction can have either sign even when AI has the same productivity multiplier in both team types. It also states exactly what a randomized factorial design identifies. No empirical interaction is estimated here.

\section{A project with two necessary components}
Each funded project runs for the same three-year horizon and must discover two distinct components to produce one validated result. A failed project has outcome zero. Conditional on effective research efforts $e_1,e_2\geq0$, the two component discoveries are independent and have probabilities $1-\exp(-r_a e_i)$, where AI access is $a\in\{0,1\}$, $r_0=1$ and $r_1=r>1$. These probabilities arise, for example, from independent Poisson search processes over the fixed horizon, with effort incorporating time and rate units.

Team composition is $m\in\{0,1\}$, where 1 denotes a prespecified mixed-discipline composition. All projects receive budget $B>0$ in common research-resource units. Mixed teams spend $c\in(0,B)$ units coordinating. For the initial benchmark AI access is supplied without deducting additional project resources. Feasible research effort therefore satisfies $e_1+e_2\leq B-cm$. A validated result earns scientific score $V_m=1+vm$, with $v\geq0$, on a scale frozen before assignment and assessed without knowledge of treatment. The higher weight represents a declared value for the kind of result that the mixed team is attempting; it is not an ex post relabeling of successful projects.

Let $Y(a,m)$ be this score, including zeros for failures. Its optimized mean is denoted $y_{am}$. Resource expenditure is common across cells: the coordination cost is already counted by reducing available productive effort. It must not be subtracted a second time from an outcome measured in scientific-score units.

\begin{theorem}[Optimal allocation between necessary discoveries]
For every treatment cell the unique optimal effort allocation is
$e_1=e_2=(B-cm)/2$, and
\[
 y_{am}=(1+vm)F(r_a(B-cm)),\qquad F(x)=(1-e^{-x/2})^2.\tag{1}
\]
\end{theorem}
\begin{proof}
Positive effort in both components yields a positive success probability, while zero effort in either yields zero. Hence an optimum has both efforts positive. The objective is strictly increasing in each effort, so the budget binds. On positive effort, maximizing the product of the two success probabilities is equivalent to maximizing the sum of their logarithms. For $q>0$, the function $\log(1-e^{-q e})$ has second derivative $-q^2e^{qe}/(e^{qe}-1)^2<0$. With equal component technologies, strict concavity and symmetry therefore give the unique equal split. Substitution proves (1). Compactness of the effort simplex also ensures the maximum is attained.
\end{proof}

\section{A complete interaction threshold}
The target factorial interaction is
\[
 \Delta=y_{11}-y_{01}-y_{10}+y_{00}.
\]
For $x>0$, define the AI improvement $D_r(x)=F(rx)-F(x)>0$. Both team types benefit from AI in this model, regardless of the sign of their interaction.

\begin{theorem}[Complementarity need not reflect a special cross-field technology]
The interaction obeys
\[
 \Delta=(1+v)D_r(B-c)-D_r(B).
\]
It is positive, zero or negative according as $v$ is greater than, equal to or less than
\[
 v^*=\frac{D_r(B)}{D_r(B-c)}-1.\tag{2}
\]
For $v=0$ and $c=B/2$, $\Delta<0$ for all sufficiently small positive $B$, but $\Delta>0$ for all sufficiently large $B$.
\end{theorem}
\begin{proof}
Subtract the four means in (1). Since $D_r(B-c)>0$, division gives the exact threshold (2), including cases in which $v^*<0$ and every admissible $v$ gives a positive interaction. Near zero,
$F(x)=x^2/4+O(x^3)$, so
\[
 D_r(x)=\frac{r^2-1}{4}x^2+O(x^3).
\]
Consequently, for $c=B/2$ and $v=0$, $\Delta/B^2\to-3(r^2-1)/16<0$ as $B\downarrow0$. At large $x$, expansion of the squares gives
$D_r(x)\sim2e^{-x/2}$ because $r>1$. Thus $D_r(B/2)/D_r(B)\sim e^{B/4}\to\infty$, proving the positive sign for sufficiently large budgets.
\end{proof}

The small-budget case reflects the need to complete both components: halving scarce productive resources sharply reduces the opportunity for AI to help. The large-budget case reflects saturation: single-discipline teams are already close to certain success, leaving less room for improvement than in the resource-constrained mixed team. The model has no additional AI advantage specific to a mixed team. A positive factorial interaction alone therefore does not identify a mechanism of cross-disciplinary translation or matching. Conversely, a negative interaction is compatible with strictly positive AI benefits in both team types.

Allowing a resource cost $k$ for AI and an AI-specific coordination saving $\eta$ changes available effort to $E_{am}=B-cm-ka+\eta am$, assumed positive in all cells. The first theorem still applies after replacing $B-cm$ by $E_{am}$, because the two-component allocation problem is unchanged. The exact interaction is then
\[
 (1+v)\{F(r(B-c-k+\eta))-F(B-c)\}
 -\{F(r(B-k))-F(B)\}.
\]
This expression separates a declared resource mechanism from the scale and saturation effects already present in (2). Without measuring those quantities, the interaction cannot separate them.

\section{A finite-cohort experiment}
Fix $N$ eligible projects before assignment, each with four potential outcomes on the same validated score scale and no cross-project spillovers. Assign exactly $n_{am}>0$ projects to each cell uniformly over all partitions with those sizes. Let $\bar Y_{am}$ be the observed cell mean and
$\widehat\Delta=\bar Y_{11}-\bar Y_{01}-\bar Y_{10}+\bar Y_{00}$.
For project $i$, write $\tau_i=Y_i(11)-Y_i(01)-Y_i(10)+Y_i(00)$. Let $S_{am}^2$ and $S_\tau^2$ be finite-population variances, with denominator $N-1$.

\begin{proposition}[What factorial randomization identifies]
The estimator is unbiased for the finite-cohort interaction $N^{-1}\sum_i\tau_i$, and its exact randomization variance is
\[
 \Var(\widehat\Delta)=\sum_{a,m}\frac{S_{am}^2}{n_{am}}-\frac{S_\tau^2}{N}.\tag{3}
\]
The final term depends on the unobserved joint potential-outcome distribution, although the interaction mean is identified without that distribution.
\end{proposition}
\begin{proof}
Every project enters cell $(a,m)$ with probability $n_{am}/N$, giving unbiasedness of each cell mean. Sampling without replacement gives variance $(1/n_{am}-1/N)S_{am}^2$. For two distinct cells $j,k$, direct use of joint assignment probability $n_jn_k/[N(N-1)]$ for two different projects gives covariance $-S_{jk}/N$, where $S_{jk}$ is the finite-population covariance of their two potential outcomes. Expanding the signed contrast variance collects the $1/n_j$ terms and subtracts $N^{-1}$ times the variance of the signed individual contrast, which is exactly (3).
\end{proof}

Randomizing individual AI accounts while researchers share tools and ideas violates the no-spillover condition. A design can instead assign entire collaboration clusters, with its target and variance recalculated at that level. Changing team membership after observing AI performance likewise changes the intervention from the fixed composition considered above. Nonpublication cannot be treated as missing success: the outcome protocol must retain failures and validate claimed discoveries comparably across cells.

\section{Completion Estimate}
56\%

This is a post hoc, heuristic estimate for the full catalogue target.
The attempt solves fixed-budget effort allocation for two necessary discovery components, characterizes either sign of AI-by-team complementarity, and derives unbiased finite-cohort factorial estimation and its variance. The full discovery target still needs validated project outcomes, measured coordination burdens and spillovers, credible disciplinary comparisons, monetary conversion where required, and an institutional allocation of expertise justified by net gains.

\section{Remaining gap and reference}
The analysis assumes independent component searches, symmetric production rates and a fixed scientific value scale. It does not estimate discovery validity, coordination burdens, project selection, or longer-run spillovers. A monetary decision additionally requires an explicit value per scientific-score unit and all real costs; maximizing the interaction is not the same as choosing the most valuable treatment cell. The broad cross-disciplinary discovery question remains unresolved.

Erik Brynjolfsson, Anton Korinek and Ajay Agrawal (2024), \emph{The Economics of Transformative AI: A Research Agenda}, Section 3.2, \url{https://digitaleconomy.stanford.edu/app/uploads/2024/11/ETAI-White-Paper.pdf}. This agenda motivates the question; the two-component technology, sign reversal and randomization calculation above are the present restricted analysis.

\end{document}
